\documentclass[10pt,a4paper]{article}
\usepackage[T1]{fontenc}
\usepackage[utf8]{inputenc}
\usepackage[brazil]{babel}
\usepackage{lmodern}
\usepackage{geometry}
\geometry{margin=1.8cm}
\usepackage{microtype}
\usepackage{graphicx}
\usepackage{booktabs}
\usepackage{tabularx}
\usepackage{array}
\usepackage{float}
\usepackage{enumitem}
\setlist{nosep,leftmargin=*}
\usepackage{xcolor}
\usepackage{listings}
\usepackage{tikz}
\usetikzlibrary{positioning,arrows.meta,fit,shapes.geometric}
\usepackage[hidelinks]{hyperref}

\definecolor{codebg}{RGB}{246,248,250}
\lstset{
  basicstyle=\ttfamily\footnotesize,
  backgroundcolor=\color{codebg},
  frame=single,
  rulecolor=\color{black!15},
  breaklines=true,
  columns=fullflexible,
  keepspaces=true,
  showstringspaces=false,
  xleftmargin=3pt,
  xrightmargin=3pt
}

\title{Centralização de Logs, Auditoria e Segurança com a Stack ELK}
\author{Sistemas Operacionais Modernos -- CEFET-MG}
\date{Prática 03}

\begin{document}
\maketitle
\vspace{-1.2em}

\section{Introdução}

Aplicações e servidores web produzem registros de acesso e de erro que permitem acompanhar requisições, investigar falhas e identificar comportamentos anormais. Quando esses registros permanecem dispersos em arquivos locais, a análise torna-se trabalhosa. A centralização permite pesquisar eventos em um único ambiente e correlacionar informações, como horário, endereço IP, recurso solicitado e código HTTP.

O objetivo desta prática é implantar uma infraestrutura de centralização e consulta de logs utilizando Elasticsearch, Filebeat e Kibana. Como serviço web, foi utilizada uma aplicação Flask publicada por meio do servidor Nginx. O Nginx registra as requisições em \texttt{access.log} e as falhas de acesso a arquivos em \texttt{error.log}; o Filebeat coleta esses arquivos e envia os eventos ao Elasticsearch; e o Kibana permite pesquisar e visualizar os registros.

\section{Arquitetura e fluxo de dados}

Os serviços são executados em containers Docker conectados à rede virtual \texttt{stackelk-network}, definida no arquivo \texttt{docker-compose.yml}. O fluxo da solução é apresentado na Figura~\ref{fig:arquitetura}.

\begin{figure}[H]
\centering
\begin{tikzpicture}[
  node distance=5mm and 7mm,
  caixa/.style={draw,rounded corners,thick,align=center,minimum height=9mm,text width=2.35cm,fill=blue!5,font=\small},
  log/.style={draw,rounded corners,thick,align=center,minimum height=9mm,text width=2.55cm,fill=orange!10,font=\small},
  seta/.style={-{Latex[length=2mm]},thick}
]
\node[caixa] (cliente) {Navegador /\texttt{gerar\_trafego.py}};
\node[caixa,right=of cliente] (nginx) {Nginx\\porta 8080};
\node[caixa,right=of nginx] (flask) {Aplicação\\Flask};
\draw[seta] (cliente)--(nginx);
\draw[seta] (nginx)--node[above,font=\scriptsize]{proxy HTTP}(flask);
\node[log,below=11mm of nginx] (logs) {Volume compartilhado\\\texttt{access.log} e \texttt{error.log}};
\node[caixa,right=of logs] (beat) {Filebeat\\coleta};
\node[caixa,right=of beat] (es) {Elasticsearch\\indexação};
\node[caixa,right=of es] (kibana) {Kibana\\Discover / Lens};
\draw[seta] (nginx)--node[left,font=\scriptsize]{grava}(logs);
\draw[seta] (logs)--(beat);
\draw[seta] (beat)--(es);
\draw[seta] (es)--(kibana);
\end{tikzpicture}
\caption{Arquitetura de coleta, armazenamento e consulta dos logs.}
\label{fig:arquitetura}
\end{figure}

O navegador e o script de teste acessam o Nginx pela porta 8080. As rotas dinâmicas são encaminhadas à aplicação Flask. O Nginx registra todas as requisições no arquivo \texttt{access.log}, em formato JSON, e as falhas de abertura de arquivos estáticos no \texttt{error.log}. Os dois arquivos ficam em um volume Docker compartilhado: Nginx possui acesso de escrita e Filebeat realiza a leitura.

O Filebeat utiliza a rede interna e o nome do serviço \texttt{elasticsearch} como destino. Cada linha do log de acesso é interpretada como um documento JSON; as linhas do log de erro são encaminhadas como mensagens textuais. O Elasticsearch indexa os eventos em índices diários \texttt{filebeat-<versão>-AAAA.MM.DD}. O Kibana consulta o padrão \texttt{filebeat-*}.

\section{Ferramentas e utilização no projeto}

\textbf{Nginx e Flask.} Flask fornece a página inicial, a rota \texttt{/sobre}, a API \texttt{/api/status} e o endpoint de saúde \texttt{/health}. Nginx atua como servidor de entrada e proxy reverso. O formato JSON do log de acesso registra \texttt{timestamp}, \texttt{client.ip}, método HTTP, URI, código de resposta, tamanho da resposta e agente do usuário. O formato estruturado facilita a extração de campos e as agregações no Elasticsearch.

\textbf{Filebeat.} O arquivo \texttt{filebeat/filebeat.yml} declara dois inputs do tipo \texttt{filestream}: um monitora \texttt{/var/log/stackelk/access.log} e utiliza o parser \texttt{ndjson}; o outro monitora \texttt{/var/log/stackelk/error.log}. Um processador converte o campo temporal do log de acesso para \texttt{@timestamp}. A saída \texttt{output.elasticsearch} envia os eventos ao serviço Elasticsearch pela porta 9200 na rede Docker.

\textbf{Elasticsearch.} Recebe os eventos do Filebeat e os armazena em índices pesquisáveis. O campo \texttt{client.ip} permite agrupar requisições por origem, enquanto \texttt{http.response.status\_code} permite filtrar respostas HTTP 404. A configuração utiliza um nó único e desabilita a segurança apenas para simplificar o laboratório local.

\textbf{Kibana.} A Data View \texttt{filebeat-*} permite explorar os eventos em Discover. A pesquisa pode ser refinada por período e por campos estruturados. No Lens, é possível construir uma tabela ou gráfico com a contagem de registros por \texttt{client.ip}, limitada às respostas 404.

\textbf{Docker Compose.} O arquivo \texttt{docker-compose.yml} define os cinco serviços, a rede virtual e os volumes persistentes. As dependências com condição \texttt{service\_healthy} evitam iniciar os componentes dependentes antes dos serviços essenciais estarem disponíveis.

\section{Execução e validação}

A implantação é realizada na raiz do projeto. Em sistemas Linux, antes de iniciar o Elasticsearch, ajusta-se o limite de memória virtual:

\begin{lstlisting}
sudo sysctl -w vm.max_map_count=262144
docker compose pull
docker compose up -d --build
docker compose ps
\end{lstlisting}

A aplicação fica disponível em \texttt{http://localhost:8080} e o Kibana em \texttt{http://localhost:5601}. A Data View é criada pelo script \texttt{scripts/criar\_data\_view.sh}.

A validação utiliza \texttt{scripts/gerar\_trafego.py}, que não exige bibliotecas adicionais. Primeiro, são geradas 30 requisições legítimas, distribuídas entre as rotas existentes:

\begin{lstlisting}
python3 scripts/gerar_trafego.py --normal-only --count 30
\end{lstlisting}

Em seguida, o script tenta acessar 250 arquivos inexistentes em \texttt{/arquivos/inexistentes/}:

\begin{lstlisting}
python3 scripts/gerar_trafego.py --attack-only --count 250 --delay 0.01
\end{lstlisting}

A rota de arquivos estáticos é atendida diretamente pelo Nginx. Para cada arquivo inexistente, o servidor registra uma resposta HTTP 404 no \texttt{access.log} e a falha de abertura no \texttt{error.log}. Assim, o teste exercita os dois inputs configurados no Filebeat.

No Kibana, em \textit{Discover}, seleciona-se a Data View \texttt{filebeat-*}, ajusta-se o intervalo de tempo para abranger a execução e aplica-se a consulta KQL:

\begin{lstlisting}
event.dataset : "nginx.access" and http.response.status_code : 404
\end{lstlisting}

As colunas \texttt{@timestamp}, \texttt{client.ip}, \texttt{http.response.status\_code} e \texttt{url.original} mostram o horário, a origem, o código e o recurso requisitado. Para examinar o arquivo de erro, utiliza-se \texttt{event.dataset : "nginx.error"}. Uma visualização Lens com a métrica \textit{Count of records} e agrupamento por \texttt{client.ip} evidencia as origens com maior quantidade de falhas.

\begin{table}[H]
\centering
\small
\caption{Cenários de validação disponibilizados no projeto.}
\begin{tabularx}{\linewidth}{>{\raggedright\arraybackslash}p{3.2cm} p{2.5cm} X}
\toprule
\textbf{Cenário} & \textbf{Volume planejado} & \textbf{Evidência esperada} \\
\midrule
Tráfego legítimo & 30 requisições & Respostas HTTP 200 nas rotas existentes, registradas no \texttt{access.log}. \\
Varredura de recursos & 250 requisições & Respostas HTTP 404, IP de origem e URIs inexistentes no Discover; falhas de abertura no \texttt{error.log}. \\
\bottomrule
\end{tabularx}
\end{table}

Os volumes da tabela correspondem aos parâmetros do script de validação. A contagem final deve ser confirmada no Kibana após a execução, pois depende de o ambiente estar ativo e de todos os eventos terem sido encaminhados pelo Filebeat.

\newpage
\section{Conclusão}

O sistema integra um serviço web, um servidor Nginx e os componentes da Stack Elastic para centralizar logs de acesso e erro. O Nginx produz registros estruturados e textuais; o Filebeat monitora os arquivos e encaminha os eventos; o Elasticsearch os indexa; e o Kibana permite pesquisar e visualizar as ocorrências. O Docker Compose torna a infraestrutura reproduzível, organizando os serviços em uma rede comum e compartilhando os arquivos por volume.

Os dois cenários automatizados permitem comparar acessos legítimos com uma sequência de tentativas a recursos inexistentes. A consulta KQL por código 404 e o agrupamento por IP permitem localizar o intervalo da anomalia e identificar a origem observada pelo servidor. A configuração foi preparada para laboratório acadêmico; em produção, seria necessário habilitar autenticação, TLS e políticas adequadas de controle de acesso e retenção.

\section*{Repositório e materiais}

O pacote do projeto contém os arquivos de aplicação, configuração do Nginx, Docker Compose, Filebeat e scripts de validação. A documentação de instalação e operação está no arquivo \texttt{README.md}.

\begin{thebibliography}{9}
\bibitem{elastic}
ELASTIC. \textit{Filebeat Reference}. Disponível em: \url{https://www.elastic.co/guide/en/beats/filebeat/current/index.html}. Acesso em: 21 nov. 2026.

\bibitem{es}
ELASTIC. \textit{Elasticsearch Guide}. Disponível em: \url{https://www.elastic.co/guide/en/elasticsearch/reference/current/index.html}. Acesso em: 21 nov. 2026.

\bibitem{kibana}
ELASTIC. \textit{Kibana Guide}. Disponível em: \url{https://www.elastic.co/guide/en/kibana/current/index.html}. Acesso em: 21 nov. 2026.

\bibitem{nginx}
NGINX. \textit{Configuring Logging}. Disponível em: \url{https://nginx.org/en/docs/http/ngx_http_log_module.html}. Acesso em: 21 nov. 2026.

\bibitem{docker}
DOCKER. \textit{Docker Compose}. Disponível em: \url{https://docs.docker.com/compose/}. Acesso em: 21 nov. 2026.
\end{thebibliography}

\end{document}
